\input{../tex-inputs/latex-leseansicht-vorspann}

\section[Arthur Schnitzler an Gustav Schwarzkopf, 2{[}3?{]}. 7. 1913]{L04508 Arthur Schnitzler an Gustav Schwarzkopf, 2{[}3?{]}. 7. 1913}
\nopagebreak\mylabel{L04508v}
\rehead{ }\normalsize\beginnumbering\briefempfaengerindex{Schwarzkopf, Gustav@\textsc{Schwarzkopf, Gustav}!zzzSchnitzler, Arthur@\emph{von Arthur Schnitzler}!1913-07-232@{23. 7. 1913}|(be}
\toendnotes[C]{\smallbreak\pagebreak[2]}
\correspDesc{Versand  durch Arthur Schnitzler am 23. 7. 1913 in Wien
\newline{}Übermittlung  am 2[3?]. 7. 1913 in Wien
\newline{}Weiterleitung  am 24. 7. 1913 in Opatija
\newline{}Erhalt  durch Gustav Schwarzkopf im Zeitraum [28. 7. 1913
                  – 5. 8. 1913?] in Wien}\toendnotes[C]{\smallbreak}
\Standort{DLA, A:Schnitzler, HS.1985.1.1897.}
\physDesc{Bildpostkarte, 229 Zeichen
\newline{}Handschrift: schwarze Tinte, deutsche Kurrent
\newline{}Versand: 1) Stempel: »\nobreak{}\oindex{Wien@\textbf{Wien}|pwk}Wien 110, 2\textcolor{gray}{3}. VII. 13, 3\nobreak{}«.   2) ursprüngliche Adresse durchgestrichen und umadressiert nach: »Wien I., Tiefer Graben
                                    17\oindex{Wien@\textbf{Wien}!I., Innere Stadt@\textbf{I., Innere Stadt}!Tiefer Graben 17@\textbf{Tiefer Graben 17}|pw}.« 3) Stempel: »\nobreak{}\oindex{Opatija@\textbf{Opatija}|pwk}Abbazia, 24. VII. 13, 7\nobreak{}«. }\toendnotes[C]{\smallbreak}
\pstart
           \noindent{}{\pb}\textcolor{gray}{\textbf{Wien, XVIII, Sternwartestr. 71\oindex{Wien@\textbf{Wien}!XVIII., Währing@\textbf{XVIII., Währing}!Sternwartestraße 71@\textbf{Sternwartestraße 71}|pw}.}}\pend
           \pstart{}{\pb}Herr Guſtav Schwarzkopf\pend{}\pstart{}Abbazia\oindex{Opatija@\textbf{Opatija}|pw},\pend{}\pstart{}Wiener Heim\oindex{Pension Wiener Heim@\textbf{Pension Wiener Heim}|pw}\pend{}{\bigskip}\vspace{1em}
\pstart
           \raggedleft{}2\textcolor{gray}{3}. 7.\pend
           \vspace{0.5em}
\pstart
           lieber Guſtav, ich weiſs nicht ob Sie \label{K_L04508-1v}\edtext{meine letzte Karte}{\lemma{\textnormal{\emph{meine letzte Karte}}}\Cendnote{\textnormal{XXXX Auszeichnungsfehler: Dokument L04223 nicht gefunden. Offenbar hatte Schwarzkopf\pwindex{Schwarzkopf, Gustav 7.\,11.\,1853 Wien – 13.\,11.\,1939 ebd.@\textsc{Schwarzkopf, Gustav} (7.\,11.\,1853 Wien – 13.\,11.\,1939 ebd.)|pwk} die Karte mit der Information zum geplanten
                  Ankunftsdatums der Familie Schnitzler\pwindex{Schnitzler, Olga 17.\,1.\,1882 Wien – 13.\,1.\,1970 Lugano@\textsc{Schnitzler, Olga} (17.\,1.\,1882 Wien – 13.\,1.\,1970 Lugano)|pwkv}\pwindex{Schnitzler, Heinrich 9.\,8.\,1902 Hinterbrühl – 12.\,7.\,1982 Wien@\textsc{Schnitzler, Heinrich} (9.\,8.\,1902 Hinterbrühl – 12.\,7.\,1982 Wien)|pwkv}\pwindex{Cappellini, Lili 13.\,9.\,1909 Wien – 26.\,7.\,1928 Venedig@\textsc{Cappellini, Lili} (13.\,9.\,1909 Wien – 26.\,7.\,1928 Venedig)|pwkv} in
                     Brijuni\oindex{Brijuni@\textbf{Brijuni}|pwk} erhalten, denn am 25. 7. 1913, dem Tag darauf, traf er
                  selbst dort ein und besuchte sie an den folgenden beiden Tagen. Die vorliegende
                  Karte erreichte ihn dagegen erst nach seiner Rückkehr nach Wien\oindex{Wien@\textbf{Wien}|pwk}, wie aus der Umadressierung zu schließen ist, also erst
                  nach dem 27. 7. 1913. Er dürfte, als die Karte in seinem Urlaubsort
                     Opatija\oindex{Opatija@\textbf{Opatija}|pwk} eintraf schon auf dem Weg nach
                     Pula\oindex{Pula@\textbf{Pula}|pwk} gewesen sein, wo er sich
                  einquartierte, um von dort nach Brijuni\oindex{Brijuni@\textbf{Brijuni}|pwk}
                  überzusetzen.}}}\label{K_L04508-1} erhielten. Für alle Fälle theile ich Ihnen mit: heute Abend fahren wir nach Brioni\oindex{Brijuni@\textbf{Brijuni}|pw} u hoffen
               Sie bald zu ſehen. Herzlichſt Ihr\pend
           \pstart \spacefill\mbox{Arthur}\pend{}\selectlanguage{ngerman}\endnumbering\briefempfaengerindex{Schwarzkopf, Gustav@\textsc{Schwarzkopf, Gustav}!zzzSchnitzler, Arthur@\emph{von Arthur Schnitzler}!1913-07-232@{23. 7. 1913}|)be}\mylabel{L04508h}
\begin{anhang}
\end{anhang}\newcommand{\dateiname}{L04508}\newcommand{\titel}{Arthur Schnitzler an Gustav Schwarzkopf, 2[3?]. 7. 1913}\newcommand{\editorInnen}{Herausgegeben von Jahnke, SelmaMüller, Martin Anton}\input{../tex-inputs/latex-leseansicht-abspann}
